%% ma: L11-B14-0025
%% lop: 11
%% chuong: 4
%% bai: 14
%% dang: 11.14.1
%% loai: DS
%% muc: [TH, TH, VD, NB]
%% dap_an: "ĐĐSS"
%% nguon: "tai-lieu-hoc-tap-mon-toan-lop-11-theo-chuong-trinh-giao-duc-pho-thong-2018.pdf (Lớp Toán Thầy Duy, Phạm Lê Duy, THPT Chu Văn An), Chương 4 Bài 5, Phần 2, Câu 5 (tr.532)"
%% kien_thuc: "Lớp 11 Bài 14 (phép chiếu song song)"
%% trang_thai: chinh-thuc
%% ghi_chu: "Bảng đáp án gốc ý c) là Đ; theo kho (kiểm bằng tọa độ) ý c) là Sai vì A và B' cùng chiếu thành B' nên ảnh suy biến thành đoạn thẳng. Cần giáo viên đối chiếu đề gốc (có thể đề gốc in sai phương chiếu hoặc tên điểm). Đã duyệt 03/11/2026."
\begin{ex}
Cho hình hộp chữ nhật $ABCD.A^{\prime}B^{\prime}C^{\prime}D^{\prime}$.
\choiceTF
{\True $ADC^{\prime}B^{\prime}$ là hình bình hành.}
{\True Hình chiếu song song của $A$ lên mặt phẳng $(BCC^{\prime}B^{\prime})$ qua phép chiếu song song theo phương của đường thẳng $C^{\prime}D$ là $B^{\prime}$.}
{Hình chiếu song song của $\triangle AB^{\prime}D^{\prime}$ trên mặt phẳng $(BCC^{\prime}B^{\prime})$ qua phép chiếu song song theo phương $C^{\prime}D$ là $\triangle BB^{\prime}C^{\prime}$.}
{Hình chiếu của hình hộp $ABCD.A^{\prime}B^{\prime}C^{\prime}D^{\prime}$ lên mặt phẳng $(BC^{\prime}D)$ theo phương $AB$ là một hình tam giác.}
\loigiai{\sloppy a) $AD\parallel B^{\prime}C^{\prime}$ và $AD=B^{\prime}C^{\prime}$ nên $ADC^{\prime}B^{\prime}$ là hình bình hành. Đúng.
b) $AB^{\prime}\parallel DC^{\prime}$ (hình bình hành $AB^{\prime}C^{\prime}D$) nên $A\mapsto B^{\prime}$. Đúng.
c) Theo phương $C^{\prime}D$: $A\mapsto B^{\prime}$, $B^{\prime}\mapsto B^{\prime}$ và $D^{\prime}\mapsto D_1$ (với hệ tọa độ $A(0;0;0)$, $B(1;0;0)$, $D(0;1;0)$, $A^{\prime}(0;0;1)$: $A\mapsto(1;0;1)=B^{\prime}$, $D^{\prime}(0;1;1)\mapsto(1;1;2)=D_1$). Hai đỉnh $A,B^{\prime}$ cùng có ảnh $B^{\prime}$ nên ảnh của $\triangle AB^{\prime}D^{\prime}$ là một đoạn thẳng, không phải $\triangle BB^{\prime}C^{\prime}$. Sai.
d) Hình chiếu của hình hộp theo phương $AB$ lên $(BC^{\prime}D)$ là hình bình hành (xem Phần 1, Câu 13), không phải tam giác. Sai.}
\end{ex}
